\chapter{Revisão Bibliográfica}
\label{cap2}

\section{Introdução ao Capítulo}
\label{sec:revbiblio_intro}

\paragraph{}Este capítulo apresenta a revisão bibliográfica que fundamenta o desenvolvimento do \textit{MinervaMesh}, organizada em três blocos: o \textbf{estado da arte} dos temas que sustentam o trabalho --- problemas térmicos e de escoamento, métodos numéricos para EDPs, sua consolidação histórica e a emergência da computação científica no navegador; os \textbf{trabalhos relacionados} do mesmo grupo de pesquisa, que estabelecem a tradição de implementação em Python científico na qual este trabalho se insere; e, por fim, o \textbf{posicionamento do trabalho}, com a lacuna identificada e a enumeração explícita das contribuições.

\section{Estado da Arte}
\label{sec:estado_da_arte}

\subsection{Problemas Térmicos e de Escoamento}
\paragraph{}A base física dos problemas tratados pela plataforma é a da transferência de calor e da mecânica dos fluidos. A condução é descrita pela Lei de Fourier e pela equação da difusão térmica --- que em regime estacionário se reduz às equações de Laplace e Poisson ---, com tratamento de referência em engenharia consolidado por Incropera et al.\ \cite{Incropera2006} e, em registro mais introdutório, por Çengel e Ghajar \cite{Cengel2015}. Quando há interação sólido-fluido, a convecção passa a governar a taxa de remoção de calor e exige o acoplamento entre as equações de energia e de escoamento: os fundamentos físicos desse acoplamento vão dos conceitos introdutórios de Fox et al.\ \cite{Fox2012} à análise de escoamentos viscosos de White \cite{White2006}. A solução numérica conjunta de calor e fluidos é desenvolvida por Maliska \cite{Maliska2004} e apresentada didaticamente por Versteeg e Malalasekera \cite{Versteeg2007}. A radiação térmica, relevante no balanço global em aplicações de alta temperatura, permanece fora do escopo deste trabalho e é considerada extensão natural para desenvolvimentos futuros.

\subsection{Métodos Numéricos para EDPs}
\paragraph{}Dois métodos de discretização concentram a literatura de interesse. O MDF aproxima derivadas por expansões de Taylor sobre malhas estruturadas e, por sua implementação direta, é frequentemente adotado em validações iniciais e na discretização temporal de problemas transientes \cite{Patankar1980}; em problemas dominados por difusão, esquemas implícitos como o de Crank-Nicolson são preferidos pela robustez numérica \cite{Versteeg2007}. O MEF, por sua vez, oferece flexibilidade geométrica por meio de malhas não estruturadas e de uma formulação variacional, característica decisiva em domínios complexos; seu desenvolvimento e fundamentação rigorosa são consolidados no tratado de referência de Zienkiewicz et al.\ \cite{Zienkiewicz2013}. A dedução da formulação de Galerkin --- que transforma a EDP na forma fraca e conduz à montagem das matrizes globais de rigidez, massa e convecção --- é apresentada em detalhe por Reddy \cite{Reddy2005}; a aplicação do método a problemas acoplados de transferência de calor e escoamento é tratada por Lewis et al.\ \cite{Lewis2004}; e textos introdutórios como o de Fish e Belytschko \cite{Fish2007} mostram que, para elementos triangulares lineares, a implementação resultante é eficiente e adequada a plataformas computacionais didáticas. O desenvolvimento matemático desses métodos, na forma empregada pela plataforma, é objeto do Capítulo~\ref{cap3}.

\subsection{Consolidação Histórica do MEF e da CFD}
\paragraph{}O desenvolvimento do MEF resulta de uma sucessão de contribuições distintas, e não de um marco único. Galerkin \cite{Galerkin1915} introduziu o método dos resíduos ponderados que leva seu nome, no qual a solução aproximada é obtida impondo a ortogonalidade do resíduo às próprias funções de base, dispensando a existência prévia de um princípio variacional. Courant \cite{Courant1943} formalizou o emprego de funções de interpolação contínuas por partes sobre subdomínios triangulares na solução de problemas de equilíbrio e vibração, antecipando o conceito de elemento. Hrennikoff \cite{Hrennikoff1941}, por sua vez, propôs o \textit{framework method}, no qual o meio contínuo elástico é discretizado por uma treliça de barras com propriedades equivalentes --- um precursor direto da ideia de discretização do domínio. A estruturação do MEF moderno aplicado à engenharia veio com Turner et al.\ \cite{Turner1956}, que formularam a análise matricial de rigidez de estruturas aeronáuticas complexas a partir de elementos triangulares e retangulares, e com Clough \cite{Clough1960}, a quem se atribui a cunhagem do próprio termo \textit{método dos elementos finitos}.

\paragraph{}No tratamento do transporte convectivo e das equações de Navier-Stokes, a literatura desenvolveu formulações estabilizadas para suprimir as oscilações espúrias da formulação de Galerkin padrão em regime de convecção dominante. Christie et al.\ \cite{Christie1976} propuseram funções de peso assimétricas (\textit{upwind}) para equações de segunda ordem com termos de primeira derivada significativos, um dos primeiros tratamentos consistentes do problema. Brooks e Hughes \cite{Brooks1982} generalizaram essa ideia na formulação SUPG (\textit{Streamline Upwind/Petrov-Galerkin}), que adiciona difusão artificial apenas na direção das linhas de corrente e tornou-se referência para escoamentos incompressíveis. Donea \cite{Donea1984} introduziu o método Taylor-Galerkin, que deriva o esquema estabilizado a partir de uma expansão de Taylor no tempo, obtendo baixa difusão numérica em problemas de advecção transiente. Löhner et al.\ \cite{Lohner1984} estenderam essa abordagem a sistemas hiperbólicos não lineares, realizando a discretização temporal antes da espacial para recuperar uma forma auto-adjunta compatível com o processo de Galerkin.

\paragraph{}A validação de solucionadores de escoamento apoia-se em casos de referência consagrados, cada um associado a um regime distinto. Ghia et al.\ \cite{Ghia1982} produziram, por método multigrid, soluções de alta resolução para a cavidade com tampa móvel (\textit{lid-driven cavity}) em números de Reynolds de até 10.000, estabelecendo o \textit{benchmark} mais difundido para esse problema; Marchi et al.\ \cite{Marchi2009} refinaram posteriormente esse padrão com soluções em malha $1024 \times 1024$ e extrapolação de Richardson. Para o degrau regressivo (\textit{backward-facing step}), Armaly et al.\ \cite{Armaly1983} forneceram medições experimentais por anemometria laser-Doppler acompanhadas de análise numérica, enquanto Thomas et al.\ \cite{Thomas1981} apresentaram uma das primeiras análises dessa mesma geometria por elementos finitos. No campo dos algoritmos, Pironneau \cite{Pironneau1982} formalizou o método de transporte-difusão baseado em características, com aplicação direta às equações de Navier-Stokes. Também se destaca a relevância das ferramentas de geração de malha na qualidade da solução, com destaque para o Gmsh em fluxos de trabalho acadêmicos e de pesquisa \cite{Geuzaine2009}.

\subsection{Computação Científica no Navegador}
\label{sec:computacao_navegador}
\paragraph{}Mais recentemente, a maturidade do \textit{WebAssembly} abriu caminho para a execução de código científico diretamente no navegador, sem instalação local. O projeto PyScript \cite{PyScript2023} disponibiliza em ambiente web o interpretador CPython e parte de sua pilha científica --- incluindo o NumPy, cujo modelo de programação vetorial é descrito por Harris et al.\ \cite{Harris2020}. Ferramentas consolidadas de pré-processamento começam a percorrer o mesmo caminho: o \texttt{triangle-wasm} \cite{TriangleWasm} demonstra a viabilidade de portar geradores de malha clássicos para \textit{WebAssembly}. No contexto de formação em engenharia, o material didático de Anjos \cite{Anjos2024} sistematiza o emprego do Python científico na solução de problemas de engenharia, evidenciando tanto o potencial pedagógico dessa pilha quanto a barreira de configuração de ambiente que este trabalho se propõe a eliminar.

\section{Trabalhos Relacionados}
\label{sec:trabalhos_relacionados}

\paragraph{}Em projetos de graduação recentes do mesmo grupo de pesquisa, observa-se forte convergência para a implementação dos métodos numéricos em Python científico aplicada a problemas térmicos e de escoamento. Em vez de agrupá-los, os trabalhos mais próximos são resumidos individualmente a seguir, pois cada um aborda um problema físico e uma estratégia de implementação distintos.

\subsection{Transferência de Calor e Condução em Componentes}

\paragraph{}Aguiar \cite{Aguiar2020} desenvolveu uma rotina em Python para o dimensionamento térmico preliminar de discos de freio, resolvendo a equação axissimétrica de condução de calor por MEF e validando o código por convergência de malha e confronto com soluções analíticas. Alves \cite{Alves2021} estendeu o problema do disco de freio para a equação tridimensional transiente do calor, também por MEF, com ênfase em oferecer alternativa de código aberto aos pacotes comerciais para comparar materiais e geometrias. Ferreira \cite{Ferreira2022} aplicou o MEF em Python à análise térmica de processadores computacionais, contrastando o comportamento térmico sob diferentes discretizações em condições críticas de operação. Pinheiro \cite{Pinheiro2022} tratou a condução de calor em componentes eletrônicos de material heterogêneo, modelando a dissipação resistiva como geração interna e resolvendo o problema bidimensional por elementos finitos.

\paragraph{}Capitanio \cite{Capitanio2023} combinou MEF (Galerkin no espaço) e MDF (no tempo) para o estudo paramétrico de dissipadores de processadores, validando a metodologia por solução manufaturada, com erro quadrático médio da ordem de $10^{-4}$ a $10^{-6}$ conforme o refino de malha. Miranda \cite{Miranda2023} investigou geometrias de canais em placas frias para arrefecimento de eletrônicos, acoplando transferência de calor e escoamento pela formulação função-corrente--vorticidade e validando contra o escoamento de Poiseuille. Teixeira \cite{Teixeira2023} resolveu a equação de Laplace tridimensional para a temperatura em um bloco de motor, empregando uma cadeia de ferramentas integralmente de código aberto (FreeCAD, Gmsh, Python/Julia) e comparando fluidos de arrefecimento. Oliveira \cite{Oliveira2025} otimizou o número de aletas de um dissipador de calor de cobre por MEF em Python, validando o solver contra solução analítica com erro relativo da ordem de $10^{-11}$.

\subsection{Escoamento e CFD pela Formulação Vorticidade-Corrente}

\paragraph{}Na linha de mecânica dos fluidos, há forte convergência para a formulação vorticidade--corrente das equações de Navier-Stokes implementada em Python. Florentino \cite{Florentino2022} aplicou essa formulação, com estabilização Taylor-Galerkin, ao arrefecimento de eletrônicos, validando o código em casos clássicos (Poiseuille, cilindro em canal) antes de simular canais aletados. Santos \cite{Santos2022} empregou a mesma formulação para investigar os campos de velocidade, vorticidade e função corrente em escoamento livre ao redor de diferentes geometrias, com a verificação apoiada nos escoamentos de Poiseuille e na cavidade com tampa móvel. Silva \cite{SilvaJessica2020} acoplou a formulação vorticidade--corrente a uma descrição arbitrária lagrangiana-euleriana (ALE) para simular as oscilações de um cilindro sob escoamento transversal, abordando a interação fluido-estrutura. Amaral \cite{Amaral2020} e Gomes \cite{GomesMatheus2021} estenderam a formulação ao transporte convectivo-difusivo de espécie química em artérias coronárias --- com ênfase, respectivamente, nas diferentes geometrias de \textit{stent} farmacológico e nas configurações de restrição de fluxo ---, evidenciando a necessidade de esquemas estabilizados em regime de convecção dominante.

\subsection{Comparação de Métodos e Fundamentação Variacional}

\paragraph{}Dois trabalhos tratam diretamente da comparação entre métodos e da fundamentação matemática do MEF. Araujo \cite{Araujo2022} comparou solução analítica, MDF e MEF em problemas bidimensionais de condução de calor de complexidade crescente --- até o caso transiente com geração de calor, sem solução fechada --- caracterizando o comportamento do erro com o refinamento de malha; é a referência mais próxima do contraste MDF\,$\times$\,MEF retomado neste trabalho. Gomes \cite{GomesYuri2025} apresentou a análise variacional da equação de Poisson em transferência de calor, estabelecendo existência e unicidade da solução fraca pelo teorema de Lax-Milgram e validando a discretização de Galerkin (MEF em Python) por solução manufaturada, com erro relativo de aproximadamente 0,46\%.

\paragraph{}Em conjunto, esses trabalhos consolidam, no mesmo grupo de pesquisa, uma tradição de implementação dos métodos numéricos em Python científico para problemas térmicos e de escoamento. Todos, contudo, pressupõem a instalação local do ambiente Python e a execução em \textit{desktop}; nenhum disponibiliza as simulações de forma interativa na web --- lacuna sobre a qual o \textit{MinervaMesh} se posiciona.

\section{Posicionamento do Trabalho e Contribuições}
\label{sec:posicionamento_contribuicoes}

\paragraph{}A revisão indica que: (i) há base teórica consolidada para problemas térmicos e de escoamento; (ii) MDF e MEF são métodos complementares, com o MEF mais adequado para geometrias complexas; e (iii) existe demanda por ferramentas didáticas acessíveis para implementação e validação desses métodos. Os trabalhos relacionados confirmam a maturidade da implementação desses métodos em Python científico dentro do próprio grupo de pesquisa, mas nenhum deles elimina a barreira de instalação e configuração de ambiente: as simulações permanecem confinadas ao \textit{desktop} de quem as desenvolveu.

\paragraph{}É exatamente nessa lacuna que este trabalho se enquadra: levar a tradição de implementação aberta em Python científico do grupo para um meio de entrega sem barreira de acesso --- o navegador ---, apoiando-se na maturidade recente do \textit{WebAssembly} descrita na Seção~\ref{sec:computacao_navegador}. Frente a esse panorama, as contribuições deste trabalho são:
\begin{itemize}
    \item o desenvolvimento da plataforma \textit{MinervaMesh}, ambiente web de código aberto para simulação numérica em Engenharia Mecânica;
    \item a execução integral das simulações no navegador, via \textit{WebAssembly} e \textit{PyScript}, sem instalação nem licenciamento;
    \item a disponibilização de código Python aberto e inspecionável em cada simulação, preservando a transparência dos algoritmos;
    \item a implementação de módulos baseados em MEF e em MDF, cobrindo os domínios térmico, estrutural e fluidodinâmico;
    \item a validação dos módulos frente a soluções analíticas e a \textit{benchmarks} clássicos da literatura;
    \item a aplicação voltada ao ensino de Engenharia Mecânica, para uso por alunos desta instituição e de outras.
\end{itemize}
